
\documentclass[10pt]{article}
\usepackage{graphicx}
\usepackage{amsmath}
\usepackage{amssymb}
\usepackage{amsfonts}
\usepackage{textcomp}
\usepackage{amsthm}

%% ---- Theorem Environments ---------------------------------
\theoremstyle{plain}
\newtheorem{theorem}{Theorem}[section]
\newtheorem{lemma}[theorem]{Lemma}
\newtheorem{corollary}[theorem]{Corollary}
\newtheorem{proposition}[theorem]{Proposition}

\theoremstyle{definition}
\newtheorem{definition}[theorem]{Definition}
\newtheorem{example}[theorem]{Example}

\theoremstyle{remark}
\newtheorem{remark}[theorem]{Remark}

\title{Equations \& Theorems for Mathematics}
\author{A. E. Mahrouss\\amlal@nekernel.org}
\date{\today}

\begin{document}

\maketitle
\vspace{-0.5em}
\vspace{0.8em}
\rule{\linewidth}{0.4pt}
\vspace{1em}

\section{Equations \& Theorems}
\begin{theorem}
    There exists:
    \begin{align*}
        \frac{\int (np) \text{ } \operatorname{dp}}{\int (px) \text{ } \operatorname{dx}} &= \frac{p^2n+2C}{px^2+2C} \in \mathbb{D}^{+}.
    \end{align*}
    That satisfies:
    \begin{align*}
        \int (px) \operatorname{dp} &= \int (np) \operatorname{dx} \frac{p^2n+2C}{px^2+2C};\\\\
        px^2+2C &= \int (np) \text{ } \operatorname{dp};\\
        \int (px) \operatorname{dp} &= px^2+2C \frac{p^2n+2C}{px^2+2C};\\
        \int (px) \operatorname{dp} &= p^2n+2C;\\\\
        \frac{p^2n+2C}{px^2+2C} &= \frac{p^2n+2C}{px^2+2C};\\
        \frac{p^2n+2C}{px^2+2C} - \frac{p^2n+2C}{px^2+2C} &= 0.
    \end{align*}
\end{theorem}

\nocite{*}

\bibliographystyle{acm}
\bibliography{A_Course_of_Pure_Mathematics, General_Theory_Dirichlet_Series, Riemann}

\end{document}